% Generated draft: requires mathematical and source review.
\documentclass[11pt]{article}
\usepackage[margin=25mm]{geometry}
\usepackage{fontspec}
\setmainfont{Latin Modern Roman}
\usepackage{amsmath,amssymb,mathtools,unicode-math}
\setmathfont{Latin Modern Math}
\usepackage{enumitem,needspace,xurl,hyperref}
\setlength{\parindent}{0pt}
\setlength{\parskip}{5pt}
\setlength{\emergencystretch}{3em}
\newcommand{\N}{\mathbb{N}}
\newcommand{\Z}{\mathbb{Z}}
\newcommand{\Q}{\mathbb{Q}}
\newcommand{\R}{\mathbb{R}}
\newcommand{\C}{\mathbb{C}}
\newcommand{\F}{\mathbb{F}}
\newcommand{\A}{\mathbb{A}}
\newcommand{\PP}{\mathbb{P}}
\newcommand{\E}{\mathbb{E}}
\newcommand{\eps}{\varepsilon}
\newcommand{\dd}{\,\mathrm{d}}
\begin{document}
\section*{UnsolvedMath ID: 30002202}
\textbf{Source problem code:} catalogue entry
\subsection{Definitions and mathematical statement}
A polyhedral fan \(\Sigma\) is a collection of polyhedral cones in a common real vector space, closed under taking faces, such that the intersection of any two cones is a face of each. The supplied record does not specify the dimension, finiteness, rationality, completeness, or purity assumptions on \(\Sigma\). Its background identifies the intended setting as a nonsimplicial toric piecewise-polynomial module, but does not give the precise admissible class of fans.

The notation \(C^0(\Sigma)\) denotes the module of compatible polynomial pieces on the cones of \(\Sigma\): polynomial pieces must agree on common faces. To make this definition algebraically precise, one must fix the coefficient ring, the rings of polynomial functions on the cones, and the ring over which freeness is being asked. None of these choices is specified in the supplied record.

More explicitly, a conditional formulation is available once these data are fixed. Let \(S\) be the chosen ring of global polynomials, let \(S_\sigma\) be the chosen ring of polynomial functions on a cone \(\sigma\), and let \(r_{\sigma,\tau}:S_\sigma\to S_\tau\) be the compatible restriction maps for faces \(\tau\subseteq\sigma\). Assume these rings carry compatible \(S\)-algebra structures. Then the compatible-piece module is
\[
C^0(\Sigma)=\left\{(f_\sigma)_{\sigma\in\Sigma}\in\prod_{\sigma\in\Sigma}S_\sigma:
 r_{\sigma,\tau}(f_\sigma)=f_\tau\ \text{whenever }\tau\text{ is a face of }\sigma\right\},
\]
with its induced \(S\)-module structure. This describes the required choices without identifying them as verified conventions of the cited report. An \(S\)-module \(M\) is free if there is a set \(I\) and an \(S\)-module isomorphism \(M\cong\bigoplus_{i\in I}S\). Whether graded freeness or finite rank is intended is also unspecified.

\textbf{Question (with unresolved source parameters).} For the intended class \(\mathcal C\) of fans and the intended polynomial-module conventions, does there exist a simple combinatorial condition \(P\) such that
\[
\forall\Sigma\in\mathcal C,\qquad
P\bigl(D(\Sigma)\bigr)
\quad\Longleftrightarrow\quad
C^0(\Sigma)\text{ is free over }S,
\]
where \(D(\Sigma)\) is the combinatorial data permitted by the intended formulation?

Here \(\mathcal C\), \(S\), the allowable data \(D(\Sigma)\), and the meaning of ``simple'' remain unresolved. In particular, the record does not say whether the condition may use only the face poset or also lattice or geometric data. The displayed equivalence is therefore a statement schema requiring source clarification, rather than a fully specified conjecture.
\subsection{Short English statement}
The problem asks for a simple combinatorial criterion that holds exactly when the piecewise-polynomial module \(C^0(\Sigma)\) is free. The supplied background places it in a nonsimplicial toric setting. The coefficient ring, base ring for freeness, admissible fans, permitted combinatorial data, and meaning of ``simple'' need clarification before the question has a fully precise formulation.
\subsection{Sources}
\begin{itemize}
\item Local UnsolvedMath registry: supplied complete record, permanent ID \(30002202\), title \emph{Combinatorial Freeness of Piecewise-Polynomial Fan Modules}, exact source problem code \texttt{catalogue entry}, in the Oberwolfach Reports Open Problems collection. No individual proposer is identified; the supplied proposed year is 2012.
\item Supplied report citation: \emph{Mini-Workshop: Cohomology Rings and Fundamental Groups of Hyperplane Arrangements, Wonderful Compactifications, and Real Toric Varieties} (2013), Oberwolfach Reports, pp.~2939--2983, DOI \texttt{10.4171/owr/2012/49}. Supplied source URL: \url{https://doi.org/10.4171/owr/2012/49}. The record describes Oberwolfach Reports as documenting mathematical workshops and their open-problem discussions.
\item Supplied literature source, labeled ``original OWR report'': \url{https://doi.org/10.4171/owr/2012/49}. The literature assessment dated \texttt{2026-08-21} labels the problem open and states that a simple combinatorial criterion for freeness of the relevant nonsimplicial toric piecewise-polynomial module remains an open problem in the cited report. This is supplied dated curation, not an independently verified literature assessment; the record itself requests independent verification of time-sensitive status and cited claims.
\item The registry supplies the statement-status label \texttt{corrected\_verified} and states that the mathematical content was checked against the official report, with notation, context, and question form normalized without changing the source-level claim. This verification claim is supplied metadata and has not been independently checked here. The report text itself is not included in the supplied evidence.
\end{itemize}
\end{document}
